\begin{textsample}{NEG Dim 5 – human – Score -9.00 – mg/mg\_ef\_linguagens\_arte\_1\_p4.txt}  \label{ex:f5_neg_014}
Essas competências e \textbf{habilidades}, \textbf{essenciais} para a formação humana, cuja arte pode desenvolver e potencializar em cada \textbf{estudante}, devem ser \textbf{trabalhadas} nas diversas linguagens do presente componente \textbf{curricular}, e que na \textbf{Base} \textbf{Nacional} Comum \textbf{Curricular} ( BNCC ) e no Currículo Referência de Minas Gerais, vêm \textbf{descritos} como unidades temáticas.
Ao \textbf{longo} do Ensino Fundamental, os \textbf{estudantes} devem expandir seu repertório e ampliar sua autonomia nas práticas artísticas, por meio da reflexão sensível, imaginativa e crítica sobre os conteúdos artísticos e seus elementos constitutivos e também sobre as experiências de pesquisa, invenção e criação.

% matched lemmas: base, curricular, descrever, essencial, estudante, habilidade, longo, nacional, trabalhar
\end{textsample}
